% Part: normal-modal-logic
% Chapter: completeness
% Section: complete-consistent-sets

\documentclass[../../../include/open-logic-section]{subfiles}

\begin{document}

\olfileid[fr]{nml}{com}{ccs}

\olsection{Ensembles \texorpdfstring{$\Sigma$}{Σ}-cohérents complets}

Supposons que $\Sigma$ soit un ensemble de !!{formula}s modales : on
peut les considérer comme les axiomes ou les principes définissant une
logique modale normale. Un ensemble $\Gamma$ est $\Sigma$-cohérent si
et seulement si $\Gamma \Proves/[\Sigma] \lfalse$, c'est-à-dire s'il
n'existe aucune !!{derivation} de~$!A_1 \lif (!A_2 \lif \cdots (!A_n \lif
\lfalse)\dots)$ à partir de $\Sigma$, où chaque $!A_i \in \Gamma$.
Nous construirons un modèle « canonique » dans lequel chaque monde est
un ensemble $\Sigma$-cohérent d'un type particulier : il est non
seulement $\Sigma$-cohérent, mais maximal en ce sens qu'il détermine la
valeur de vérité de toute !!{formula} modale. Pour chaque $!A$, ou bien
$!A \in \Gamma$, ou bien $\lnot !A \in \Gamma$ :

\begin{defn}
  Un ensemble $\Gamma$ est \emph{$\Sigma$-cohérent complet} si et
  seulement s'il est $\Sigma$-cohérent et si, pour tout~$!A$, ou bien
  $!A \in \Gamma$, ou bien $\lnot !A \in \Gamma$.
\end{defn}

Les ensembles $\Sigma$-cohérents complets, tels que $\Gamma$, ont plusieurs
propriétés utiles. D'abord, ils sont déductivement clos : si $\Gamma
\Proves[\Sigma] !A$, alors $!A \in \Gamma$. En particulier, toute
instance par substitution d'une !!{formula}~$!A \in \Sigma$ appartient aussi à~$\Gamma$. De plus,
l'appartenance à $\Gamma$ reflète les conditions de vérité des
connecteurs propositionnels. Ce sera important lorsque nous définirons
le « modèle canonique ».

\begin{prop}\ollabel{prop:ccs-properties}
  Supposons que $\Gamma$ soit $\Sigma$-cohérent complet. Alors :
  \begin{enumerate}
  \item \ollabel{prop:ccs-closed}%
    $\Gamma$ est déductivement clos dans~$\Sigma$.
  \item \ollabel{prop:ccs-sigma}%
    $\Sigma \subseteq \Gamma$.
  \tagitem{prvFalse}{\ollabel{prop:ccs-lfalse}%
    $\lfalse \notin \Gamma$}{}
  \tagitem{prvTrue}{\ollabel{prop:ccs-ltrue}%
    $\ltrue \in \Gamma$}{}
  \item \ollabel{prop:ccs-lnot}%
    $\lnot!A \in \Gamma$ si et seulement si $!A \notin
    \Gamma$.
  \tagitem{prvAnd}{\ollabel{prop:ccs-land}%
    $!A \land !B \in \Gamma$ si et seulement si $!A \in \Gamma$ et $!B \in \Gamma$}{}
  \tagitem{prvOr}{\ollabel{prop:ccs-lor}%
    $!A \lor !B \in \Gamma$ si et seulement si $!A \in \Gamma$ ou $!B \in \Gamma$}{}
  \tagitem{prvIf}{\ollabel{prop:ccs-lif}%
    $!A \lif !B \in \Gamma$ si et seulement si $!A \notin \Gamma$ ou $!B \in \Gamma$}{}
  \tagitem{prvIff}{\ollabel{prop:ccs-liff}%
    $!A \liff !B \in \Gamma$ si et seulement si l'on a soit $!A \in \Gamma$ et $!B \in
    \Gamma$, soit $!A \notin \Gamma$ et $!B \notin \Gamma$}{}
  \end{enumerate}
\end{prop}

\begin{proof}
  \begin{enumerate}
  \item Supposons que $\Gamma \Proves[\Sigma] !A$, mais que $!A \notin
    \Gamma$. Puisque $\Gamma$ est $\Sigma$-cohérent complet, on a
    $\lnot!A \in \Gamma$. L'ensemble $\Gamma$ serait alors incohérent,
    car $!A, \lnot !A \Proves[\Sigma] \lfalse$.

  \item Si $!A \in \Sigma$, alors $\Gamma \Proves[\Sigma] !A$ ; donc $!A
    \in \Gamma$ par clôture déductive, c'est-à-dire d'après le
    point~\olref{prop:ccs-closed}.

  \tagitem{prvFalse}{Si $\lfalse \in \Gamma$, alors $\Gamma
    \Proves[\Sigma] \lfalse$ ; donc $\Gamma$ serait
    $\Sigma$-incohérent.}{}

  \tagitem{prvTrue}{On a $\Gamma \Proves[\Sigma] \ltrue$ ; donc $\ltrue \in
    \Gamma$ par clôture déductive, c'est-à-dire d'après le
    point~\olref{prop:ccs-closed}.}{}

  \item Si $\lnot!A \in \Gamma$, alors la cohérence impose $!A \notin
    \Gamma$ ; réciproquement, si $!A \notin \Gamma$, alors
    $\lnot!A \in \Gamma$, puisque $\Gamma$ est $\Sigma$-cohérent complet.

  \tagitem{prvAnd}{\iftag{probAnd}{Exercice.}{Supposons que $!A \land !B \in
      \Gamma$. Puisque $(!A \land !B) \lif !A$ est une instance de
      tautologie, $!A \in \Gamma$ par clôture déductive, c'est-à-dire
      d'après le point~\olref{prop:ccs-closed}. De même, $!B \in \Gamma$.
      Réciproquement, supposons que $!A \in \Gamma$ et $!B \in
      \Gamma$. La clôture déductive entraîne alors que $(!A \land !B) \in
      \Gamma$, puisque $!A \lif (!B \lif (!A \land !B))$ est une
      instance de tautologie.}}{}

  \tagitem{prvOr}{\iftag{probOr}{Exercice.}{Supposons que $!A \lor !B \in
      \Gamma$, tandis que $!A \notin \Gamma$ et $!B \notin \Gamma$.
      Puisque $\Gamma$ est $\Sigma$-cohérent complet, $\lnot!A \in \Gamma$
      et $\lnot !B \in \Gamma$. Dès lors, $\lnot(!A \lor !B) \in \Gamma$,
      car $\lnot !A \lif (\lnot !B \lif \lnot (!A \lor !B))$ est une
      instance de tautologie. L'ensemble $\Gamma$ serait donc
      $\Sigma$-incohérent, contradiction. Réciproquement, si $!A \in \Gamma$ ou $!B \in \Gamma$, les instances tautologiques $!A \lif (!A \lor !B)$ et $!B \lif (!A \lor !B)$ donnent $!A \lor !B \in \Gamma$ par clôture déductive.}}{}

  \tagitem{prvIf}{\iftag{probIf}{Exercice.}{Supposons que $!A \lif !B \in
      \Gamma$ et $!A \in \Gamma$ ; alors $\Gamma \Proves[\Sigma] !B$,
      d'où $!B \in \Gamma$ par clôture déductive. Réciproquement, si $!A
      \lif !B \notin \Gamma$, alors, puisque $\Gamma$ est
      $\Sigma$-cohérent complet, $\lnot (!A \lif !B) \in \Gamma$. Comme
      $\lnot(!A \lif !B) \lif !A$ est une instance de tautologie, $!A \in
      \Gamma$ par clôture déductive. Comme $\lnot(!A \lif !B) \lif
      \lnot !B$ est une instance de tautologie, $\lnot !B \in
      \Gamma$. Donc $!B \notin \Gamma$, puisque $\Gamma$ est
      $\Sigma$-cohérent.}}{}

  \tagitem{prvIff}{\iftag{probIff}{Exercice.}{Supposons que $!A \liff !B \in
      \Gamma$. Si $!A \in \Gamma$, alors $!B \in \Gamma$, car $(!A
      \liff !B) \lif (!A \lif !B)$ est une instance de
      tautologie. De même, si $!B \in \Gamma$, alors $!A \in \Gamma$.
      Ainsi, ou bien $!A \in \Gamma$ et $!B \in \Gamma$, ou bien
      $!A \notin \Gamma$ et $!B \notin \Gamma$.

      Réciproquement, supposons que $!A \liff !B \notin \Gamma$. Puisque $\Gamma$
      est $\Sigma$-cohérent complet, $\lnot (!A \liff !B) \in
      \Gamma$. Comme $\lnot(!A \liff !B) \lif (!A \lif \lnot !B)$ est une
      instance de tautologie, si $!A \in \Gamma$, alors $\lnot !B \in
      \Gamma$ et, puisque $\Gamma$ est $\Sigma$-cohérent, $!B \notin
      \Gamma$. De même, si $!B \in \Gamma$, alors $!A \notin \Gamma$.
      Ainsi, les deux formules ne peuvent appartenir à~$\Gamma$. Si elles en étaient toutes deux absentes, $\lnot!A \in \Gamma$ et la tautologie
      $\lnot(!A \liff !B) \lif (\lnot!A \lif !B)$ donneraient $!B \in \Gamma$, contradiction.}}{}
  \end{enumerate}
\end{proof}

\begin{probtag}{probAnd,probOr,probIf,probIff}
  Compléter la démonstration de la \olref[nml][com][ccs]{prop:ccs-properties}.
\end{probtag}
\begin{explain}
La complétude impose l'appartenance d'au moins une des formules
$!A$ et $\lnot!A$. La cohérence interdit qu'elles appartiennent
toutes deux à~$\Gamma$. La clôture déductive permet alors
d'utiliser les instances de tautologies pour établir les lois
d'appartenance des autres connecteurs.
\end{explain}
\end{document}
